% ==============================================================================
% Secção 4: Impacto do Projeto
% ==============================================================================

\subsection{Impacto a longo prazo}
\begin{campotexto}
A concretização das soluções propostas para os carros do futuro produz impactos significativos em múltiplas dimensões, em linha com o enquadramento PESTEL considerado no projeto. No plano \textbf{científico e tecnológico}, impulsiona a integração entre armazenamento avançado, controlo preditivo por IA e amortecimento regenerativo, estimulando investigação sobre viabilidade industrial e autonomia real. Na dimensão \textbf{económica}, as baterias de sódio, baseadas num recurso abundante, e a maturação do estado sólido favorecem a redução de custos, enquanto o aproveitamento das vibrações melhora a eficiência energética. No âmbito \textbf{social}, reforça a confiança na mobilidade elétrica ao elevar a segurança --- eliminando eletrólitos inflamáveis e aumentando a estabilidade --- e o conforto dos utilizadores. Por fim, no domínio \textbf{ambiental}, reduz a dependência de matérias-primas críticas, diminui perdas energéticas e promove a reciclagem e o descarte responsável das baterias.
\end{campotexto}

\subsection{Potencial de inovação}
\begin{campotexto}
Face às soluções automóveis existentes, as \textbf{vantagens inovadoras} deste projeto residem na integração coordenada de tecnologias emergentes. Em comparação com as baterias de lítio (LFP/NMC), o estado sólido elimina eletrólitos inflamáveis, reduz a degradação térmica e aumenta a densidade energética, enquanto as baterias de sódio substituem o lítio por um elemento abundante e mais económico. Paralelamente, a conjugação de suspensões ativas, sensores com IA e amortecedores regenerativos supera as suspensões passivas de calibração fixa e a regeneração restrita à travagem, garantindo atuação preditiva, maior estabilidade, conforto e aproveitamento das vibrações.
No \textbf{impacto no mercado}, estas vantagens reforçam a competitividade dos veículos elétricos ao aliar segurança e eficiência a menor custo, acelerando a adoção comercial após a superação dos encargos iniciais de fabrico e da complexidade de integração industrial.
\end{campotexto}

\subsection{Comunicação e divulgação}
\begin{campotexto}
A estratégia de valorização dos resultados articula a disseminação académica com a aproximação ao setor empresarial. Na vertente de \textbf{valorização científica}, prevê-se a publicação de artigos em revistas e conferências dedicadas à mobilidade elétrica, armazenamento de energia e controlo veicular, partilhando o \textit{benchmarking} tecnológico, a análise PESTEL e a arquitetura integrada proposta para os carros do futuro, além de apresentações em póster em encontros científicos. Na vertente de \textbf{valorização económica}, aposta-se na participação em feiras tecnológicas, eventos automóveis e demonstrações públicas de modelos conceptuais ou protótipos futuros, visando captar parceiros estratégicos. Esta interação permitirá transferir conhecimento para os fabricantes, validar a viabilidade industrial das soluções e avaliar o seu impacto real na autonomia e eficiência global.
\end{campotexto}
